\section{Large-Momentum Effective Theory}
\label{sec:lamet}
%---------------------------------------------------------------------
As has been explained in \sec{lc}, Feynman's partons
were motivated from describing the structure of a bound state
travelling at large momentum $P$. On the other hand, in QCD
factorizations, they appear as effective degrees of freedoms
arising in infinite momentum limit disregarding UV divergences.
Reconciling these two pictures results in large-momentum
effective theory (LaMET) for the parton structure of hadrons.

In this section, we start by considering the structure
of the proton at finite momentum. We define the
ordinary momentum distributions of the constituents,
and trying to illustrate their dependence on
the proton momentum. We demonstrate that the large-$P$ momentum-dependence
follow a RGE, similar to the
well-known RGE for partons. In \sec{lamet-fact},
we show that momentum distributions at large $P$,
are related to PDFs through a matching between different orders of
$P\to \infty$ and UV cut-off limits. This matching process has a standard EFT
explanation: Parton physics or observables
can be obtained from an effective theory in which
$P\ll \Lambda_{\rm UV}$ are calculated non-perturbatively in the so-called ${\cal P}$ space~\cite{Messiah:1979eg},
after ``integrating out'' degrees of freedom between $P$ and $\infty$  (or ${\cal Q}=1-{\cal P}$ space)
through perturbation theory. Therefore, the LaMET approach to partons
is in some sense similar to lattice QCD  as an EFT approach for continuum field theories,
in which all active degrees of freedom (${\cal P}$ space) are bounded by $|k| \le \pi/a$,
where $a$ is lattice spacing, whereas those at $|k| \ge \pi/a$ (${\cal Q}$ space) are taken into account through
perturbative coefficients and higher dimensional operators.

In \sec{lamet-pp}, we outline the formalism
of LaMET for a general parton observable. The method
can in principle be used also to calculate any LF
correlations in terms of large momentum external states (see in particular the application to
soft function in \sec{tmd}).
The strategy is also applicable for the components of the LF wave functions.
Thus, LaMET offers a practical and systematic way to carry out the program of
LFQ. Instead of working with the LF coordinates directly, one uses the
instant form of dynamics and large momentum or boost factor $\gamma$ as a
regulator for the LF divergences. In a certain sense, the
quantization using tilted light-cone coordinates~\cite{Lenz:1991sa} is
similar to the spirit of the LaMET approach.

At present, the only systematic approach to solve non-perturbative QCD
is lattice field theory~\cite{Wilson:1974sk}. Therefore,
a practical implementation of LaMET can be done
through lattice calculations.  It can also be done
with other bound-state methods using Euclidean approaches, such as the
instanton liquid model~\cite{Schafer:1996wv}. While LFQ may provide an attractive physical picture for the proton,
the Euclidean equal-time formulation is more practical for carrying out the
calculations, and LaMET serves to bridge them.

%------------------------------------------------------------------------------
\subsection{Structure of the Proton at Finite Momentum}
\label{sec:lamet-frame}
%-----------------------------------------------------------------------------
In relativistic theories, the internal structure of
a composite system is frame-dependent (we always refer to the total momentum eigenstates),
and we are interested in the properties of the proton
at a momentum much larger than its rest mass.

We start from the quark momentum density in a fast-moving proton, assuming that it moves in the $z$-direction. A
straightforward definition is
\begin{align}
	N_P(\vec{k})=\sum_\sigma \langle P|b^\dagger_\sigma (\vec{k}) b_\sigma
	(\vec{k})|P\rangle/\langle P|P\rangle \ ,
\end{align}
where the quark helicity, color, and other implicit indices are summed over.
This equation should be compared with the parton density in Eq. (\ref{eq:partonfock}).
To make it gauge invariant, it is convenient to consider
the definition from a coordinate-space correlator,
\begin{align}
	N_{P,W}(\vec{k})=\int {d^3\xi\over {(2\pi)^3}} \ e^{-i\vec{k}\cdot \vec{\xi}}
	\langle P|\overline{\psi}(0)\gamma^0W(0,\vec{\xi})\psi(\vec{\xi})|P\rangle\,,
\end{align}
where the Dirac matrix $\gamma^0$ ensures that it is a number density. Clearly,
it is a static quantity without time-dependence and can be calculated
in Euclidean field theories, in contrast to Eq. (\ref{eq:LCqPDF}) for partons.
The gauge invariance is ensured by the Wilson line $W(0,\vec{\xi})$
between the quark fields separated by $\vec{\xi}$, which is defined in the fundamental representation of the color SU(3) group. There are infinitely many
choices for the Wilson line, generating infinitely many momentum densities.
For example, one can choose a straight-line link between $0$ and $\vec{\xi}$.
One can also let the Wilson line run from the fields along the $z$-direction for
a long distance (if not infinity) before joining them together along
the transverse direction (a staple).

For its obvious connection to the PDFs, we consider a transverse-momentum
integrated, longitudinal-momentum distribution,
\begin{align}
	N_P(k^z)&=\int  d^2\vec{k}_\perp\ N_{P,W}(\vec{k}) \nonumber \\
	&=\int {dz\over{2\pi}} e^{-ik^z z} \langle P|\overline{\psi}(0)\gamma^0 W(0,z)\psi(z)|P\rangle,
	\label{eq:momdis1}
\end{align}
where we ignore the question of convergence at large $|\vec{k}_\perp| $.
Now the gauge-link $W(0,z)$ is naturally taken as a straight-line,
\begin{align}
	W(0,z) &= \exp(-i\int^z_0 A^{z'}(z') dz') \nonumber \\
	&= \exp(i\int^0_\infty A^{z'}(z') dz')\exp(i\int^\infty_z A^{z'}(z') dz')\nonumber \\
	&= W^\dagger(\infty,0) W(\infty,z)\,,
\end{align}
where in the second line we have split the gauge link into two, going from $z$ to the infinity
and coming back from the infinity to zero. We can define a ``gauge-invariant''
quark field
\begin{align}
	\Psi(\vec{\xi}) = W(\infty,\vec{\xi})\psi(\vec{\xi}) \ ,
\end{align}
and the above density becomes,
\begin{align}
	N_P(k^z)= \int {dz\over {2\pi}} e^{-ik^z z}  \langle P|\overline{\Psi}(0)\gamma^0\Psi(z)|P\rangle \ ,
	\label{eq:momdis}
\end{align}
where again we have not considered UV divergences. The momentum
distribution defined above has been called {\it quasi-PDF}, but it is really
a physical momentum distribution in a proton of momentum $P$.

In the rest frame of the proton, $N_{{P=0}}(k^z)$ is symmetric in positive and negative $ k^z$,
probably peaks around $k^z=0$ and decays away as $k^z\to \pm\infty$.
Due to the perturbative QCD effects, it decays algebraically at large $k^z$, instead of exponentially.
Because of this property, the high moments of the distribution, $\int dk^z (k^z)^n N_0(k^z)$ with $n>0$,
have the standard QFT UV divergences.

As $P$ becomes non-zero and large, the peak $N_P(k_z)$ will be around $\alpha P^z$,
where $\alpha$ is a constant of order one. The density
at negative $k^z$ becomes smaller, but not zero. This is due to the so-called
backward-moving particles from the
large momentum kick in perturbation theory. For the same reason, the density at $k^z>P^z$ is not zero
either.

$N_P(k^z)$ has a renormalization scale dependence because the quark fields must be renormalized.
One can choose DR and modified minimal subtraction ($\overline{\rm MS}$)
scheme. Any other regularization scheme can be converted into this one perturbatively.
For $z\ne 0$, the only renormalization necessary is the quark wave
function (with anomalous dimension $\gamma_{F}$) in the $A^z=0$ gauge, because the linear divergence associated with the gauge link vanishes in the $\MS$ scheme. More extensive discussions on the renormalization
issue, particularly about non-perturbative renormalization,
will be made in the following section.

As an example showing how the parton momentum density depends on $P$, we depict in Fig.~\ref{fig:thooftwf} the quark wave function amplitude of a
meson in the 't Hooft model (1+1 dimensional QCD with $N_c\to\infty$)~\cite{tHooft:1974pnl} , the square of which yields the quark momentum density. In this model, a meson of momentum $P^\mu$ can be built as
\begin{align}
	|P_n^\mu\rangle & =
	\int \frac{dk}{2\pi|P|}
	\big[M(k-P,k)\phi_n^+(k,P)  \nonumber \\
	&\qquad\qquad + M^\dagger(k,k-P)\phi_n^-(k,P)\big]|0\rangle\,,
\end{align}
where $M(p,k) = \sum_i d^i_{-p}b_k^i/\sqrt{N_c}$, and $M^\dagger(p,k)
= \sum_i b_k^{i\dagger} d_{-p}^{i\dagger}/\sqrt{N_c}$ are annihilation
and creation operators for quark-antiquark pairs. The corresponding wave function amplitudes,
$\phi_n^+(k,P)$ and $\phi_n^-(k,P)$, satisfy a pair of equations
first derived in~\cite{Bars:1977ud}.

\begin{figure}[htb]	
	\centering
	\includegraphics[width=0.95\linewidth]{images/intro-WvFn_LC_fig04.pdf}
	\caption{Wave function amplitudes of a meson in the 't Hooft model at different external momenta~\cite{Jia:2017uul}.}
	\label{fig:thooftwf}
\end{figure}

The meson bound state defined above has a well-defined large-momentum limit.
The wave functions can be expanded in
$1/P$, with the corrections starting from $(1/P)^2$.
The momenta of the constituents, $k$ and $P-k$, scale in this
limit. When plotted as a function of $x=k/P$, the change in the
wave function with the magnitude of the momentum is shown in Fig.~\ref{fig:thooftwf}.

%------------------------------------------------------------------------------
\subsection{Momentum Renormalization Group}
\label{sec:lamet-rg}
%------------------------------------------------------------------------------
In this subsection, we consider how to calculate the external momentum $P$
dependence of physical observables discussed in
the previous subsection.
Clearly, the dependence is related to the boost properties of
the operators under consideration, namely their commutation relations
with the boost generators, {$\hat{K}^i$}.
We argue that in the large momentum limit, one
has a {\it momentum RGE} which
is a differential equation relating
properties of the system at different momenta.
Momentum RGE will be, in the end, related to
the renormalization properties of the observables on
the LF.

Consider a generic operator $\hat O$, and its
matrix element in a state with momentum $P$,
\begin{align}
	O(P) = \langle P|\hat O|P\rangle \ .
\end{align}
We calculate the momentum dependence by writing $|P\rangle
= \exp(-i\omega(P) \hat K)|P=0\rangle$, where $\hat K$ is
the boost operator along the momentum direction and $\omega$ is
a boost parameter depending on $P$.
Taking a derivative with respect to the {boost} parameter gives
\begin{align}
	\frac{d O(P)}{dP} = i\frac{d\omega(P)}{dP}\langle P|[\hat O,\hat K]|P\rangle \ .
	\label{eq:momden}
\end{align}
The {r.h.s.} of the equation depends on the commutator
$[\hat O,\hat K]$, i.e., the boost properties of the operator.
For a scalar operator, the commutation relation vanishes, and $O(P)$
is frame independent.
For a vector operator, the commutation relation resembles that of an energy-momentum four-vector,
and the result is the standard Lorentz transformation of a four-vector.
For nonlocal operators, the commutation
relation requires the elementary formula,
\begin{align}
	[J^{\mu\nu}, \phi_i(x)] = i\left[l^{\mu\nu}\delta_{ij}+S^{\mu\nu}_{ij}\right]\phi_j(x)\ ,
\end{align}
where $l^{\mu\nu}\!=\!-i (x^{\mu}\partial^{\nu}\!-\!x^{\nu}\partial^{\mu})$ is the OAM operator and $S^{\mu\nu}$ is the intrinsic spin matrix.
Thus one of the fields is now { $\phi_i(t=\sinh\omega z, 0, 0, \cosh\omega z)$}
which generates a time-dependent correlation function.

In the large-momentum limit, because of asymptotic freedom,
the $P$-dependence is calculable in perturbation theory, and Eq. (\ref{eq:momden})
simplifies. One obtains the momentum or boost RGE~\cite{Ji:2014gla},
\begin{align}
	\frac{d O(P)}{dP} &= \lim_{\Delta P\to 0}[O(P+\Delta P) - O(P)]/\Delta P \\
	&\xrightarrow{P\gg M} C(\alpha_s(P)){\otimes} O(P) + {\cal O} (M^2/P^2) \ .
\end{align}
where $C(\alpha_s(P))$ is a perturbative expansion in the strong coupling $\alpha_s$. The symbol ``$\otimes$'' can be a simple multiplication or certain form of convolution, depending on the observable $O(P)$ studied.
The proof of the above equation is non-trivial, and it can be
analyzed on a case-by-case basis. There can be mixings among a set of independent operators with the same quantum numbers.
The momentum RGEs are very similar to those for scale transformation or that for the
coarse graining of a Hamiltonian. That the two are connected in some cases may
be traced to Lorentz symmetry.

\begin{figure}[t]	
	\centering
	\begin{subfigure}{.15\textwidth}
		\centering
		\includegraphics[width=\linewidth]{images/intro-quasi_quark_fig05a.pdf}
		\caption{}
		\label{fig:qoneloop1}
	\end{subfigure}
	\begin{subfigure}{.15\textwidth}
		\centering
		\includegraphics[width=\linewidth]{images/intro-quasi_quark_fig05b.pdf}
		\caption{}
		\label{fig:qoneloop2a}
	\end{subfigure}
	\begin{subfigure}{.15\textwidth}
		\centering
		\includegraphics[width=\linewidth]{images/intro-quasi_quark_fig05c.pdf}
		\caption{}
		\label{fig:qoneloop2b}
	\end{subfigure}
	\begin{subfigure}{.15\textwidth}
		\centering
		\includegraphics[width=\linewidth]{images/intro-quasi_quark_fig05d.pdf}
		\caption{}
		\label{fig:qoneloop3a}
	\end{subfigure}
	\begin{subfigure}{.15\textwidth}
		\centering
		\includegraphics[width=\linewidth]{images/intro-quasi_quark_fig05e.pdf}
		\caption{}
		\label{fig:qoneloop3b}
	\end{subfigure}
	\begin{subfigure}{.15\textwidth}
		\centering
		\includegraphics[width=\linewidth]{images/intro-quasi_quark_fig05f.pdf}
		\caption{}
		\label{fig:qoneloop4}
	\end{subfigure}
	\caption{One-loop diagrams for the quasi-PDF in a free quark state in the Feynman gauge. The conjugate diagrams of (b), (c), (e), (f) do contribute but are not shown here.}
	\label{fig:qoneloop}
\end{figure}

As an example of the momentum RGE,
we calculate the quark momentum distribution
in a perturbative quark state using Eq.~(\ref{eq:momdis}).
Since it is gauge invariant, we can calculate it
in any gauge, for example, the Feynman gauge. The one-loop diagrams in QCD are shown in Fig.~\ref{fig:qoneloop}. There are two sources of UV
divergences, one is the logarithmic divergences from the vertex and self-energy diagrams, and the other
is the linear divergence in the self-energy of the Wilson line.
For the moment, we will use transverse momentum cut-off, $\Lambda_{\rm UV}$,  as the
UV regulator.  Using $y=k^z/P^z$, the one-loop result reads
for a large momentum quark~\cite{Xiong:2013bka},
\begin{align}\label{unpolvertex}
	& \tilde q^{(1)}(y,P^z,\Lambda_{\rm UV})=\frac{\alpha_s C_F}{2\pi} \nonumber
	\\
	& \times \left\{ \begin{array} {ll} \frac{1+y^2}{1-y}\ln \frac{y}{y-1}+1
		+\frac{\Lambda_{\rm UV}}{(1-y)^2P^z}\ , & y>1 \\
		\frac{1+y^2}{1-y}\ln \frac{(P^z)^2}{m^2}+\frac{1+y^2}{1-y}\ln\frac{4y}{1-y} \\
		\, \, \, \, -\frac{4y}{1-y}+1+\frac{\Lambda_{\rm UV}}{(1-y)^2P^z}\ , & 0<y<1 \\ \frac{1+y^2}{1-y}\ln \frac{y-1}{y}-1+\frac{\Lambda_{\rm UV}}{(1-y)^2P^z}\ , & y<0 \end{array} \right.
\end{align}
where we have ignored all power-suppressed contributions and keep the leading $P^z$ dependence only.
There is an additional contribution of the form $\delta Z_1(\Lambda_{\rm UV}/P^z) \delta(y-1)$.

The above result has several interesting features:
\begin{itemize}
	\item{The distribution does not vanish outside {$[0,1]$}.
		The radiative gluon can carry a large negative momentum fraction,
		resulting in a recoiling quark carrying larger momentum than the parent quark,
		and thus $y>1$. The same gluon can also carry a momentum larger than $P^z$,
		making the active quark have $y<0$.}
	
	\item{While the above effect is easy to understand perturbatively, it is
		surprising that a scaling contribution remains outside [0,1] in the IMF.
		As the proton travels faster,
		one might think any constituent has a momentum $k^z$ positive from Lorentz
		transformation. However, the order of limits matters because no matter
		how large the parent-quark momentum is, there are always quarks with much larger momentum, i.e.,
		$k^z\gg P^z \gg \Lambda_{\rm QCD}$. In this sense, Feynman's parton model does not describe
		the exact properties of the momentum distribution
		in a large-momentum nucleon.}
	
	\item{The contribution outside $[0,1]$ at one-loop is entirely perturbative because of the
		absence of any infrared (IR) divergence. This is no longer true at two-loop level, but
		the contribution depends only on the same one-loop IR physics in $[0,1]$. }
	
	\item{The distribution for $y$ in [0,1] has a term depending on $\ln P^z$. This dependence
		reflects that the quark substructure is resolved as a function of $P^z$, an interesting
		feature of boost.
		This dependence is perturbative in the sense that the derivative is
		IR safe,
		\begin{align}
			& P^z \frac{d \tilde q(y,P^z,\Lambda_{\rm UV})}{dP^z} \nonumber \\
			&= \frac{\alpha_s C_F}{\pi} \left[\left(\frac{1+y^2}{1-y}\right)_+ -\frac{3}{2}\delta(1-y) \right] \,.
		\end{align}
		Apart from the {$\delta$-function term}, the {r.h.s.} is similar to the one-loop
		quark splitting function in DGLAP evolution~\cite{Dokshitzer:1977sg,Gribov:1972ri,Altarelli:1977zs}. Therefore one might suspect that the momentum
		dependence is closely related to the familiar renormalization scale evolution in the
		PDFs. In fact, the physics is just the other way around: {\it It is the hadron-momentum dependence
			of the physical momentum distribution that generates
			the DGLAP evolution in the infinite-momentum limit.} One can derive an all-order momentum RGE for the momentum distribution function. Momentum RGE also provides a method
		to sum over the large logarithms of the momentum.}
	
	\item{There is a singularity at $y=1$. This singularity is generated
		from soft-gluon radiation. Fortunately, this singularity combined
		with the virtual contribution yields a finite result.}
	
	\item{There is a linear divergence in the cut-off regulator,
		leading to $\Lambda_{\rm UV}/P^z$ term, which is absent in DR. Thus, to keep $1/(P^z)^2$
		power counting, it is important
		to work in a renormalization scheme where this term does not exist.}
	
\end{itemize}

We can also move on to study the hadron momentum RGEs of other
structural properties considered in the previous subsection.
In particular, the RGE for TMD distributions will lead to
the familiar rapidity RGE in the literature. We reserve these discussions to \sec{tmd}.

%------------------------------------------------------------------------------
\subsection{Effective Field Theory Matching to PDFs}
\label{sec:lamet-fact}
%------------------------------------------------------------------------------
As we have seen in the previous subsection,
the momentum distributions of the constituents (now
called quasi-PDFs in the literature) in a proton at large $P$
are different from the PDFs or LF distributions in many ways.
In particular, the momentum fraction $y$ in a physical momentum distribution
is not limited to [0,1] due to backward moving particles,
which is the case even in the $P\to\infty$ limit. In fact, the
infinite-momentum limit is not analytical due to the presence of $\ln P$.

However, partons are effective objects arising
from a different limit $\Lambda_{\rm UV}\ll P\to \infty $.
There is also an important computational advantage in taking the naive limit $P\gg \Lambda_{\rm UV}$ in
perturbative calculations: Feynman integrals have one fewer four-momentum.
Therefore, this limit of QFTs serves as a
reference system where the structure of the bound states
is manifestly {independent of the hadron momentum}, and is similar to
scale-invariant critical points at which second-order phase transitions occur
in condensed matter systems. However, the theory in the naive IMF limit
introduces additional UV divergences.

Therefore, the partons in QCD are very similar
to the infinitely heavy quarks in HQET~\cite{Manohar:2000dt}. In certain
QCD systems, heavy quarks such as the bottom quark are present, and their
masses are much larger than the typical QCD scale $\Lambda_{\rm QCD}$.
In this case, one might study the dependence on the heavy quark mass
by expanding around $m_Q=\infty$. This expansion will generally
produce a power series in $1/m_Q$. However, the limits
of taking $\Lambda_{\rm UV}\to\infty$ and infinite heavy-quark
mass {limits} are not interchangeable, due to the presence
of the large logarithms $\ln m_Q$. In an EFT approach,
one takes the $m_Q\to\infty$ limit first, this will result in
a new theory with different UV behavior, but without the heavy-quark
mass, and symmetries among very different heavy-quark
systems become manifest. The renormalization
of the extra UV divergences yields a RGE
which can be used to resum large quark-mass logarithms.

Therefore, the momentum distribution at large-$P$ differs
from the parton distributions only in the order
of limits, their IR non-perturbative physics is
the same. In asymptotically free theories such as QCD, differences (or discontinuities) in taking the limits of $P\gg \Lambda_{\rm UV}$ and $\Lambda_{\rm UV} \gg P\to\infty$ are
perturbatively calculable, as only the high-momentum modes matter.
The differences are called {\it matching coefficients}.
Therefore, one is able to write down a power expansion for the momentum-dependent distribution (quasi-PDF) in terms of the PDF~\cite{Xiong:2013bka,Ma:2014jla,Ma:2017pxb,Izubuchi:2018srq},
\begin{align}\label{eq:fact_0}
	\tilde q(y, P^z, \mu)\! &=\! \int_{-1}^{1} \frac{dx}{|x|} C\left(\frac{y}{x}, \frac{\mu}{xP^z}\right) q(x,\mu)
	\nonumber \\
	&\qquad\qquad + {\cal O}\bigg(\frac{\Lambda_{\text{QCD}}^2}{(yP^z)^2},
	\frac{\Lambda_{\text{QCD}}^2}{((1-y)P^z)^2}
	\bigg)\,,
\end{align}
where the power correction is suppressed by the parton momentum $yP^z$ and  the spectator momentum $(1-y)P^z$~\cite{Ji:2020brr}.
This expansion {may be} also called
a factorization formula, as the quasi-PDF contains all the IR physics in the PDF, and $C$ involves only UV physics.
As we will discuss extensively in
the next section, this factorization formula is true to all orders in perturbation theory.
The above relation allows us to calculate the
LF parton physics from the momentum distribution at large
$P$. Since the expansion parameter is $\Lambda_{\rm QCD}^2/(yP^z)^2$ and $\Lambda_{\rm QCD}^2/((1-y)P^z)^2$ ,
for intermediate $y$ one might not need very large $P^z$ to neglect
the power corrections.

The above relation between the two quantities
has a simple explanation in terms of the Lorentz boost:
Consider the spatial correlation along $z$ shown
in Fig. \ref{fig:Boost} in a large momentum state.
It can be seen as approaching the LF one in the rest frame of the proton.
In other words, we are using a near-LF correlation to approximate a LF
correlation. Accordingly, we can invert the above equation recursively to express the PDF in terms of
quasi-PDF with their differences being taken care of
through the perturbative matching $\tilde C$ and power corrections,
\begin{align}\label{eq:fact_1}
	q(x, \mu)\! &=\! \int_{-\infty}^{\infty} \frac{dy}{|y|} \tilde C\left(\frac{x}{y}, \frac{\mu}{yP^z}\right) \tilde q(y,P^z,\mu)
	\nonumber \\
	&\qquad\qquad + {\cal O}\bigg(\frac{\Lambda_{\text{QCD}}^2}{(xP^z)^2},
	\frac{\Lambda_{\text{QCD}}^2}{((1-x)P^z)^2}
	\bigg)\,.
\end{align}
The above equation has an EFT interpretation:
The parton physics is calculated in an effective field theory
with physical momentum scale from $0$ to $P$, whereas the physics
from degrees of freedom from $P$ to $\infty$ can be integrated out to
generate the perturbative coefficients $\tilde C$ and
the high-order terms in $1/(P^z)^2$. In constrast to HQET, the full QCD degrees of freedom are used in LaMET calculations. In other words, the effective Lagrangian of LaMET
is the standard QCD one, while the large momentum $P$ for expansions
appears only in the external states.

\begin{figure}[htb]	
	\centering
	\includegraphics[width=0.6\linewidth]{images/intro-Boost_fig06.pdf}
	\caption{The line segment in the $z$-direction in the frame of a large-momentum hadron.
		Through Lorentz boost, it is equivalent to a line segment of length $\sim \gamma z$
		close to the light-one in the hadron state of zero momentum. Here $\gamma z/\sqrt{2}$ is the length of projection of the boosted line segment to the light-cone direction $n$. Thus, we call the dimensionless
		variable $\lambda=zP^z\sim \gamma zM$ as the quasi light-cone distance.}
	\label{fig:Boost}
\end{figure}

%----------------------------------------------------------------------------
\subsection{Recipe for Parton Physics in LaMET}
\label{sec:lamet-pp}
%------------------------------------------------------------------------------

The principle of LaMET is to simulate the time-dependence of parton observables through external states at large momentum.
Thus, we can generalize the discussions in the previous subsection
to any type of physical observables for the large momentum
proton, which will be generally called quasi-parton observables.
Examples will be given in the later sections including
transverse-momentum dependent distributions and LF wave functions.

Consider any Euclidean quasi-observable $O$ which depends
on a large hadron momentum $P^z$
and UV cut-off $\Lambda_{\rm UV}\gg P^z$.
Using asymptotic freedom, we can systematically
expand the $P^z$ dependence,
\begin{align}
	O\Big(P^z,\Lambda_{\rm UV}\Big) = Z\Big({P^z\over\Lambda_{\rm UV}},
	{P^z\over \mu}\Big){\otimes}o(\mu) + {\cal O}\Big({\Lambda_{\rm QCD}^2\over (P^z)^2}\Big)+ ... \ ,
	\label{eq:Expansion}
\end{align}
where $\otimes$ refers to a convolution if appropriate, and $Z$ factorizes all the perturbative
dependence on $P^z$ and does not contain any IR divergence.
The quantity $o(\mu)$ is defined in a theory with $P^z\rightarrow \infty$, exactly
as in Feynman's parton model. In fact, $o(\mu)$ is a LF
correlation containing all the IR collinear (and soft)
singularities. The important point of the expansion is that it
may converge at moderately large $P^z$, say a few GeV, allowing
access to quantities needed for very large $P^z$ (a few TeV).
One can also use the large boost-factor $\gamma=P^z/M$ as the expansion
parameter $1/\gamma$.

Momentum dependence of the quasi-observables can be
studied through momentum RGEs. Defining the anomalous dimension through
\begin{align}
	\gamma_P(\alpha_s) = \frac{1}{Z}\frac{\partial Z}{\partial \ln P^z} \ ,
\end{align}
it follows{ that}
\begin{align}
	\frac{\partial O(P^z)}{\partial \ln P^z} = \gamma_P(\alpha_s) \otimes O(P^z) \,,
\end{align}
up to power corrections. One can resum large logarithms involving $P^z$ using the above equation.

When taking $P^z\rightarrow \infty$ first
in $O(P^z)$ before a UV regularization is imposed, one recovers from $\hat O$ the
light-cone operator $\hat o$, by construction. On the other hand, the physical matrix
element is calculated at a large $P^z$, with UV regularization such as the lattice cut-off
imposed first. Thus the difference between the matrix elements of $\hat o$ and $\hat O$ is a matter of
the order of limits. This is the standard set-up for an EFT. The different limits do
not change the IR physics. In fact, the factorization in terms of Feynman diagrams
can be proved order by order as in the renormalization program, as discussed in the
following section.

The parton physics can be calculated more directly by reversing
Eq. (\ref{eq:Expansion}) to produce an EFT expansion,
\begin{align}
	o(\mu)= \tilde Z\Big({P^z\over\Lambda_{\rm UV}},
	{P^z\over \mu}\Big){\otimes} O\Big(P^z,\Lambda_{\rm UV}\Big) + {\cal O}\Big({\Lambda_{\rm QCD}^2\over (P^z)^2}\Big)+ ... \ .
	\label{eq:Expansion1}
\end{align}
Thus, to compute any parton observable defined by an operator made of LF dynamical fields, $\hat o$, one constructs a time-independent version $\hat O$ which, under an infinite
Lorentz boost, approaches $\hat o$. Then, one calculates the matrix element
of $\hat O$ in a hadron with large momentum $P^z$ using
whatever approach (lattice QCD is an obvious choice for the time-independent operator $\hat O$)
and then uses Eq. (\ref{eq:Expansion1}) to systematically approximate the parton {observable}.
Usually the matrix element of $\hat O$ depends on $P^z$ as well as all the lattice UV artifacts.
In principle, the latter does not affect the EFT expansion and will be cancelled by the matching coefficient $\tilde Z$ and higher order terms in the expansion. However,
in practical applications such as the quasi-PDF calculations, a nonperturbative renormalization is still necessary to remove all the power divergences to ensure a continuum limit.

%------------------------------------------------------------------------------
\subsection{Universality}
\label{sec:lamet-univ}
%------------------------------------------------------------------------------
LaMET provides a framework to systematically
compute partonic observables on the LF from the properties of
a large-momentum proton.
However, the relationship is not one-to-one. There can be
infinitely many possible Euclidean operators in the large-momentum proton
which generate the same LF observable. This is because
the large-momentum physical states have built-in
collinear (as well as soft) parton modes, and upon acting on
a Euclidean operator, they help to {\it project out}
the leading LF physics. All operators projecting out
the same LF physics form a {\it universality class}.
Accordingly, in the operator formulation
for parton physics such as SCET, one uses
LF operators to project out parton physics off the external
states of any momentum, including $P=0$.

The concepts such as universality class have been used
in critical phenomena in condensed matter physics,
where systems with different microscopic Hamiltonians
can have the same scaling properties near their critical
points. Critical phenomena correspond to the IR fixed points of the scale transformation,
and are dominated by physics at long-distance scales.
In the present case, parton physics arises from the
infinite-momentum limit, $P=\infty$,
which is a UV fixed point of the momentum RGEs. It
is the longitudinal short distance (or large momentum)
physics that is relevant at the fixed point.
However, the short distance here does not mean everything
is perturbative. The part that is non-perturbative
characterizes the partonic structure of the proton. The critical region
near $P=\infty$ acts as a filter to select only the physics
that is relevant, so universality classes emerge.

In the case of unpolarized PDFs, the initial proposal in LaMET
starts from the matrix element of the following operator~\cite{Ji:2013dva},
\begin{align}
	O_1(z)=  \overline{\psi}(0)\gamma^zW(0,z)\psi(z) \ .
	\label{eq:eucbi1}
\end{align}
However, one can equally start from~\cite{Radyushkin:2016hsy,Constantinou:2017sej},
\begin{align}
	O_2(z)= \overline{\psi}(0)\gamma^0W(0,z)\psi(z) \ ,
	\label{eq:eucbi2}
\end{align}
and the leading contributions in the large-momentum expansion are the same. One can
also consider any linear combination of the two.
In~\cite{Jia:2018qee}, the calculations have been done with
these two operators in the 't Hooft model,
and the results have been compared at different
hadron momenta. For lattice simulations,
an important issue is about the operator mixing, which
depends on specific choices of the operators
in the universality class.

Another example of Euclidean operators for PDFs
is the current-current correlators
in a pure space separation,
\begin{align}
	O_3(z)=J^\mu(0)J^\nu(z) \ ,
\end{align}
where $J^\mu$ is, for example, an electromagnetic current.
This type of correlator was first considered in~\cite{Braun:2007wv,Bali:2017gfr}
for calculating pion DA, and recently
has been suggested to calculate PDFs with generalized
bilocal ``currents''~\cite{Ma:2017pxb}.
When the matrix elements are calculated in the large-momentum states,
$O(z)$ falls into the same universality class as the operators in Eqs. (\ref{eq:eucbi1}) and
(\ref{eq:eucbi2}). Instead of using light quarks as the intermediate propagator
in $O(z)$, one can have a number of other choices for LaMET applications, including
scalars~\cite{Aglietti:1998ur,Abada:2001if} and
heavy-quarks~\cite{Detmold:2005gg}. One can also similarly
work with quark bilinear operators in any physical gauge which become the light-front one in
the large momentum limit~\cite{Gupta:1993vp}.

Another important example of universality class is the gluon helicity
contribution to the spin of the proton, as we will
discuss in detail in \sec{spin}. The gluon
spin operator $\vec{E}\times \vec{A}$ is gauge-dependent. However, in physical gauges where the transverse degrees of
freedom are dynamical, its matrix element is the same in
the large-momentum limit. Therefore,
one can potentially choose different gauges to perform
calculations at finite momentum on lattice, such as Coulomb gauge $\vec{\nabla}\cdot \vec{A}=0$,
axial gauge $A^z=0$ or temporal gauge $A^0=0$.  Different gauge choices
will have different UV properties ($\ln P$) and hence different
matching conditions. However, the IR part of the matrix element
is the same~\cite{Hatta:2013gta}.

At a practical level, it is very useful to find which
operator has the fastest convergence in the LaMET expansion.
The current-current correlators use the light-quark propagator to simulate
the light-like Wilson line (sometimes called light-ray). The quasi-PDF approach not only starts from
a quantity with clear physical meaning (a momentum distribution),
but also generates the needed Wilson line simply by rotating a
space-like one, shown in Fig. \ref{fig:Boost}. Thus, it is plausible
that the quasi-PDF will provide mathematically the fastest large-$P$
convergence than the other choices.

%------------------------------------------------------------------------------
